\section*{Hoofdstuk \ref{ch:g9:alternator} --- Van centrale tot huis: de wisselstroomgenerator}

\begin{solution}{exo:g9:alternator:1}
Alleen verandering induceert; snellere beweging, een sterkere
\omterm{def:g1:magnets:magnet}{magneet} en meer windingen geven een grotere
\omterm{def:g8:voltage-voltmeter:voltage}{spanning}; een omgekeerde beweging keert de stroom om. Een
\omterm{def:g1:magnets:magnet}{magneet} in rust induceert, hoe sterk ook, precies niets.
\end{solution}

\begin{solution}{exo:g9:alternator:2}
Elke halve slag van de rotor brengt een pool naar de spoel toe (die de
ene kant op induceert) en voert hem dan weg (die de andere kant op
induceert): de \omterm{def:g8:voltage-voltmeter:voltage}{spanning} moet bij elke halve omwenteling van
teken wisselen --- het wisselen zit in het draaien ingebouwd.
\end{solution}

\begin{solution}{exo:g9:alternator:3}
Vijftig slagen per \omterm{def:g2:measuring-time:second}{seconde} (voor de eenvoudigste machine, die met twee polen).
Elke centrale moet precies hetzelfde draairitme aanhouden --- het hele
netwerk draait in de pas en deelt één \omterm{def:g8:sound-pitch-loudness:frequency}{frequentie}.
\end{solution}

\begin{solution}{exo:g9:alternator:4}
De turbine en de generator. Alleen de duwer van de bladen
verandert: de stoom van een vlam, de stoom van een atoom, een vallend
meer, of de \omterm{def:g2:air-around-us:wind}{wind}.
\end{solution}

\begin{solution}{exo:g9:alternator:5}
De opgeslagen hoogte van het meer $\to$ de beweging van vallend water
$\to$ de rotatie van de turbine $\to$ de elektriciteit van de
generator $\to$ transformator en net $\to$ lamp. De Zon vulde het
meer: \omterm{def:g6:water-states:changes}{verdamping} tilde de zee op, regen bezorgde haar in de bergen.
\end{solution}

\begin{solution}{exo:g9:alternator:6}
De \omterm{def:g5:everyday-energy:energy}{energie} van het \omterm{def:g3:first-electric-circuit:bulb}{lampje} moet ergens vandaan komen, en de enige
voorraad op een fiets is de fietser: de dynamo zet beenwerk om in
\omterm{def:g1:light-and-shadows:source}{licht}, dus verschijnen de \omterm{def:g9:kinetic-energy-safety:joule}{joule} van het \omterm{def:g3:first-electric-circuit:bulb}{lampje} als
extra tegenstand. \omterm{def:g5:everyday-energy:energy}{Energie} wordt omgezet, nooit getoverd --- de
oudste wet van de cursus.
\end{solution}

\begin{solution}{exo:g9:alternator:7}
\omterm{ex:g9:alternator:chains}{Hernieuwbaar}: \omterm{def:g2:air-around-us:wind}{wind}, het meer (de regen vult het bij),
zonlicht. Niet: kolen, uranium --- op menselijke tijdschaal voorgoed
verbruikt. De sortering vraagt: vult de natuur de voorraad bij terwijl
wij haar leegtrekken?
\end{solution}

\begin{solution}{exo:g9:alternator:8}
$P = U I$ laat één \omterm{def:g9:electric-power-energy:power}{vermogen} reizen als hoge
\omterm{def:g8:voltage-voltmeter:voltage}{spanning} met een kleine stroom; de
\omterm{prop:g8:electric-current-effects:three}{warmtewerking} laat het verlies in de draden met de stroom
groeien. Kleine stroom, koele draden, gespaarde
\omterm{def:g5:everyday-energy:energy}{energie} --- dus verscheept het net hoog en stapt het bij de
voordeur af.
\end{solution}

\begin{solution}{exo:g9:alternator:9}
Langzaam draaien: een kleine geïnduceerde \omterm{def:g8:voltage-voltmeter:voltage}{spanning} (regel twee)
--- een zwak \omterm{def:g3:first-electric-circuit:bulb}{lampje} --- bij weinig kringlopen per \omterm{def:g2:measuring-time:second}{seconde}: zichtbaar
geflikker, onder de samensmeltdrempel van het oog. Snel draaien: een
grotere \omterm{def:g8:voltage-voltmeter:voltage}{spanning} --- helder --- en een
\omterm{def:g8:sound-pitch-loudness:frequency}{frequentie} die het oog gladstrijkt. Twee knoppen, één
slinger.
\end{solution}

\begin{solution}{exo:g9:alternator:10}
Bij \qty{230}{V}: $I = \num{1e8} \div 230 \approx \qty{4.3e5}{A}$ ---
vierhonderdduizend \omterm{def:g8:intensity-ammeter:intensity}{ampère}: geen kabel op aarde. Bij
\qty{400000}{V}: $I = \num{1e8} \div \num{4e5} = \qty{250}{A}$ ---
gewone transportkabel. De zaak van de transformator, gesloten.
\end{solution}

\begin{solution}{exo:g9:alternator:11}
Beide ketens luiden: voorraad $\to$ \omterm{def:g5:heat-and-insulation:heat}{warmte} $\to$ kokend water
$\to$ beweging van stoom $\to$ rotatie $\to$ elektriciteit ---
ingewikkelde waterkokers die turbines voeden. De ketens verschillen in
precies één schakel: wat de \omterm{def:g5:heat-and-insulation:heat}{warmte} maakt --- de scheikunde van
brandende kolen, of het splijten van uraniumkernen.
\end{solution}

\begin{solution}{exo:g9:alternator:12}
Stroomopwaarts: de \qty{230}{V} van de muur kwam door de
wijktransformator (\omterm{def:g8:voltage-voltmeter:voltage}{spanning} afgestapt), langs masten op honderden
kilovolt (transportvorm: \omterm{def:g9:alternating-current:dcac}{wisselstroom} op hoge
\omterm{def:g8:voltage-voltmeter:voltage}{spanning}, die een beetje \omterm{def:g5:heat-and-insulation:heat}{warmte} lekt), vanaf een
generator (rotatie naar elektriciteit, door inductie),
rondgedraaid door een turbine (de beweging van stoom of water), gevoed
door een voorraad --- zeg het opgetilde meer van de dam, opgetild door
de \omterm{def:g6:water-states:changes}{verdamping} van de Zon. Vandaar de zin van klas drie: het
\omterm{def:g1:light-and-shadows:source}{licht} in je kamer is zonneschijn --- opgeslagen, gevallen,
rondgedraaid en verscheept.
\end{solution}

\begin{solution}{pb:g9:alternator:1}
\textbf{1.} De hoogte van het meer $\to$ de razende beweging van het
water $\to$ de rotatie van de turbine $\to$ de geïnduceerde
wisselspanning van de generator $\to$ het railstel
van de centrale.
\textbf{2.} De rotormagneet draait; de spoelen staan vast. Zonder
beweging verandert het magnetisme door de spoelen niet --- en zonder
verandering geeft inductie helemaal niets.
\textbf{3.} De \omterm{def:g8:sound-pitch-loudness:frequency}{frequentie} van de \omterm{def:g8:voltage-voltmeter:voltage}{spanning}
kopieert de rotatie: verkeerde \omterm{def:g7:motion-average-speed:formula}{snelheid}, verkeerde
\omterm{def:g8:sound-pitch-loudness:frequency}{frequentie} --- en een machine die uit de pas loopt met de
\qty{50}{Hz} van het net vecht tegen elke andere
generator op het netwerk in plaats van ze te helpen.
\textbf{4.} Een sterkere rotormagneet induceert inderdaad meer ---
en de turbine zwoegt onmiddellijk harder om hem tegen de stuggere
magnetische tegenwerking in te laten draaien: meer elektriciteit eruit
vraagt meer water erdoor, \omterm{def:g9:kinetic-energy-safety:joule}{joule} voor \omterm{def:g9:kinetic-energy-safety:joule}{joule}. Het
meer betaalt, niet de \omterm{def:g1:magnets:magnet}{magneet}.
\textbf{5.} $I = \num{2e7} \div \num{2e5} = \qty{100}{A}$.
\textbf{6.} $I = \num{2e7} \div 230 \approx \qty{8.7e4}{A}$ ---
zevenentachtigduizend \omterm{def:g8:intensity-ammeter:intensity}{ampère}: onmogelijke kabel. De taak van de
transformator is precies deze ruil --- \omterm{def:g8:voltage-voltmeter:voltage}{spanning} omhoog, stroom
omlaag, \omterm{def:g9:electric-power-energy:power}{vermogen} onveranderd.
\textbf{7.} Meer water door de schuiven. Te weinig gevoed vertragen de
belaste machines als een fietser op een helling --- en zakt de
\omterm{def:g8:sound-pitch-loudness:frequency}{frequentie} van het net onder de vijftig: alle klokken en
machines op het netwerk voelen een hongerig net.
\textbf{8.} Alleen het eerste vakje --- de \omterm{def:g2:air-around-us:wind}{wind}, de duw op de
bladen. Turbine, generator en transformator stonden er volkomen
gezond bij, windstil.
\textbf{9.} Het net is één groot \omterm{def:g5:series-parallel-circuits:parallel}{parallel} netwerk: vele wegen verbinden
elke centrale met elke stad, dus kan het aandeel van een doorbroken
lijn via de overgebleven \omterm{def:g5:series-parallel-circuits:parallel}{takken} stromen --- parallelle
onafhankelijkheid, op nationale schaal.
\textbf{10.} De omgeleide, overbelaste lijnen lieten de geleverde
\omterm{def:g8:voltage-voltmeter:voltage}{spanning} een minuut lang onder nominaal zakken: de lampen van de
loods dimden en haar \omterm{ex:g6:circuits-and-safety:motor}{motoren} draaiden lui rond --- ondervoed,
zoals elk apparaat onder zijn \omterm{def:g8:voltage-voltmeter:nominal}{nominale spanning} wel moet zijn.
\textbf{11.} ``Elke \omterm{def:g9:electric-power-energy:power}{watt} van vannacht kwam uit vallend
meerwater; de Zon verdampte de zee en de regen droeg haar naar onze
bergen. Wij laten hier zonneschijn draaien --- opgeslagen in een meer,
via de hemel.''
\textbf{12.} Bijvoorbeeld: ``Beweging langs magneten maakte elke
\omterm{def:g8:voltage-voltmeter:voltage}{volt} van de nacht --- en om zeven uur vroegen tienduizend
waterkokers het meer een tikje sneller te vallen.''
\end{solution}
